\begin{frame}{실험: 노이즈 특징을 붙여가면?}
  2차원 \kb{moons}(반달 두 개) 데이터, $k=7$ 고정,
  \textbf{정보 하나도 담지 않는} 잡음 특징을 0$\to$5$\to$20$\to$50개
  붙여가며 검증 정확도 측정:
  \small
  \begin{tabular}{@{}lc@{}}
    \toprule
    잡음 특징 수 (총 차원) & 검증 정확도 ($k=7$) \\
    \midrule
    0 (2차원)  & \textbf{0.867} \\
    5 (7차원)  & 0.813 \\
    20 (22차원) & 0.680 \\
    50 (52차원) & \textbf{0.507} \\
    \bottomrule
  \end{tabular}
  \vskip1em
  \begin{block}{해석}
    52차원이면 \textbf{0.507} --- 동전 던지기(0.500)와 사실상
    무차이. 원래 들어 있던 \textbf{유용한 2차원 신호는 그대로}인데,
    \textbf{아무것도 알려주지 않는 축 50개}를 추가했을 뿐
    $\Rightarrow$ 순수하게 차원의 저주만으로 정확도가 \textbf{무너짐}.
    ``특징을 더 많이 모은 것''이 오히려 \textbf{해로운 축}을 추가한 셈.
  \end{block}
  \kq{차원의 저주는 ``차원의 \textbf{숫자}''가 아니라,
  거리에 \textbf{노이즈 축이} 얼마나 섞여 있는가}에 달려 있다.
\end{frame}
